\section{Summary}\label{sec:summary}

\noindent A simulation study of low-energy neutron hit detection in the nHCal was performed. The study tested whether a nonuniform longitudinal sampling structure can improve the response relative to the uniform nHCal baseline. The geometry was divided into three longitudinal segments, with independent absorber and scintillator thicknesses in each segment, giving a six-dimensional design space. Candidate geometries were simulated with DD4hep/DDsim using a neutron G4GPS source over 0.1--2.5 GeV/c and evaluated using detection efficiency, average fired-tile count, and average fired-layer count.

\noindent A surrogate-guided search was used to explore this design space. The initial 150-point Latin-hypercube scan broadly covered the allowed thickness bounds, and later proposal batches used a LightGBM surrogate to rank new quasi-randomly sampled candidates. The best confirmed geometry was already present in the initial scan, while later surrogate-guided proposals produced comparable candidates but did not consistently outperform it after higher-statistics reruns. This indicates that the scanned response landscape is shallow and that near-top candidates are sensitive to Monte Carlo fluctuations.

\noindent In the final matched comparison, the optimum improved the neutron detection efficiency from $0.53172 \pm 0.00331$ to $0.54184 \pm 0.00276$, corresponding to a relative gain of about 2\%. The average fired-tile count increased from $1.09068 \pm 0.00897$ to $1.13452 \pm 0.00660$, and the average fired-layer count increased from $0.97710 \pm 0.00786$ to $1.01348 \pm 0.00553$. The energy-binned comparison showed that the largest baseline--optimum separation occurs in the intermediate neutron kinetic-energy region, while the threshold scan showed that the ordering is stable near the nominal 0.5-MIP threshold but changes at very low threshold.

\noindent These results demonstrate that the surrogate-guided workflow can be used to explore and optimize the nHCal geometry design space, and that geometry-dependent changes in neutron hit-detection efficiency are observable in simulation. Within the tested bounds, the best confirmed gain is modest. The current direction can identify small improvements, but higher-statistics scans and additional systematic studies would be needed to determine whether stronger improvements exist elsewhere in the design space.
